% Part: sets-functions-relations
% Chapter: arithmetization
% Section: cuts
%
\documentclass[../../../include/open-logic-section]{subfiles}

\begin{document}

\olfileid[pt]{sfr}{arith}{cuts}
\olsection{De $\Rat$ para $\Real$}

Em essência, a propriedade de completude mostra que qualquer ponto $\alpha$ da
reta real divide essa reta em duas metades de modo perfeito: numa, $\alpha$ é o
supremo; noutra, $\alpha$ é o ínfimo. Para \emph{construir} os números reais a
partir dos racionais, Dedekind sugeriu que pensássemos nos reais simplesmente
como os \emph{cortes} que particionam os racionais. Isto é, identificamos
$\sqrt{2}$ com o \emph{corte} que separa os racionais $< \sqrt{2}$ dos racionais
$> \sqrt{2}$.

Vamos pôr isto em ordem. Se cortarmos os racionais em duas metades, podemos
identificar de modo único a partição que fizemos considerando apenas a metade
\emph{inferior}. Sendo precisos, apresentamos a seguinte definição:

\begin{defn}[Corte] 
Um \emph{corte} $\alpha$ é qualquer segmento inicial próprio não vazio dos
racionais sem máximo. Isto é, $\alpha$ é um corte sse:
\begin{enumerate}
	\item \emph{não vazio, próprio}: $\emptyset \neq \alpha \subsetneq \Rat$
	\item \emph{inicial}: para todos $p,q \in \Rat$: se $p < q \in \alpha$ então $p \in \alpha$
	\item \emph{sem máximo}: para todo $p \in \alpha$ existe $q \in \alpha$ tal que $p < q$ 
\end{enumerate} 
Então $\Real$ é o conjunto dos cortes. 
\end{defn}

Podemos pois dizer que $\sqrt{2} = \Setabs{p \in \Rat}{p^2 < 2\text{
ou }p < 0}$. Claro que há que verificar que isto \emph{é} de fato um corte,
mas relegamo-lo a \olref[check]{sec}.

Como antes, tendo definido certas entidades, precisamos a seguir de definir
funções e relações básicas sobre elas. Começamos por uma fácil:
\begin{align*}
	\alpha \leq \beta \text{ sse }\alpha \subseteq \beta
\end{align*}
Esta definição de ordem permite \emph{enunciar} o resultado central, a saber que
o conjunto dos cortes tem a propriedade de completude. Por extenso, o enunciado
tem esta forma. Se $S$ é um conjunto não vazio de cortes com majorante, então
$S$ tem supremo. Com mais pormenor: existe um corte, $\lambda$, que é majorante
de $S$, isto é, $(\forall \alpha \in S)\alpha \subseteq
\lambda$, e $\lambda$ é o menor tal corte, isto é, $(\forall \beta \in \Real)((\forall \alpha \in S)\alpha \subseteq \beta \lif \lambda \subseteq \beta)$. Segue-se a prova do resultado:

\begin{thm}\ollabel{realcompleteness}
O conjunto dos cortes tem a propriedade de completude. 
\end{thm}

\begin{proof}
Seja $S$ um conjunto não vazio de cortes com majorante. Seja $\lambda
= \bigcup S$. 
%\Setabs{p \in \Rat}{(\exists \alpha \in S)p \in \alpha}$$
Começamos por alegar que $\lambda$ é um corte:
\begin{enumerate}
\item Como $S$ tem majorante, pelo menos um corte pertence a $S$, logo
$\emptyset \neq \lambda$. Como $S$ é um conjunto de cortes, $\lambda
\subseteq \Rat$. Como $S$ tem majorante, algum $p \in \Rat$ está ausente de
todo o corte $\alpha \in S$. Logo $p\notin \lambda$, e portanto
$\lambda \subsetneq \Rat$.
\item Suponhamos $p < q \in \lambda$. Então existe algum $\alpha \in S$
tal que $q \in \alpha$. Como $\alpha$ é um corte, $p \in \alpha$. Logo
$p \in \lambda$.
\item Suponhamos $p \in \lambda$. Então existe algum $\alpha \in S$ tal
que $p \in \alpha$. Como $\alpha$ é um corte, existe algum $q \in
\alpha$ tal que $p < q$. Logo $q \in \lambda$. 
\end{enumerate}
Isto prova a alegação. Além disso, claramente $(\forall \alpha \in S)\alpha
\subseteq \bigcup S = \lambda$, isto é, $\lambda$ é majorante de $S$. Suponhamos
agora que $\beta \in \Real$ também é majorante, isto é, $(\forall \alpha \in S)\alpha \subseteq \beta$. Para qualquer $p \in \Rat$, se $p \in \lambda$, então existe $\alpha \in S$ tal que $p \in \alpha$, logo $p \in \beta$. Generalizando, $\lambda \subseteq \beta$. Logo $\lambda$ é o supremo de $S$.
\end{proof}

Temos pois um conjunto de entidades que satisfaz a propriedade de completude. Uma
forma de o dizer é: não há ``lacunas'' nos nossos cortes. (Assim: tomar novos
``cortes'' de reais, em vez de racionais, não produziria objetos novos
interessantes.)

Segue-se a definição de algumas operações sobre os reais. Começamos por embutir os
racionais nos reais estipulando $p_\Real =
\Setabs{q \in \Rat}{q < p}$ para cada $p \in \Rat$. Definimos então:
\begin{align*}
	\alpha + \beta &= \Setabs{p + q}{p \in \alpha \land q \in \beta}\\
%	\alpha - \beta &\defis \Setabs{p - q}{p \in \alpha \land q \in \Rat \setminus \beta}\\
	\alpha \times \beta &= 
	\Setabs{p \times q}{0 \leq p \in \alpha \land 0 \leq q \in \beta} \cup 0^\mathbb{R} & \text{se }\alpha, \beta \geq 0_\Real
%	\alpha \div \beta &\defis 
%	\Setabs{\nicefrac{p}{q}}{p \in \alpha \land q \in \Rat \setminus \beta}  & \text{if }\alpha \geq 0_\Real, \beta > 0_\Real
\end{align*}
Para tratar os outros casos da multiplicação, definamos primeiro:
\begin{align*}
	-\alpha &= \Setabs{p - q}{p < 0 \land q \notin \alpha}
\end{align*}
e estipulemos:
\begin{align*}
	\alpha \times \beta &\defis 
	\begin{cases}
		\mathord{-}\alpha \times \mathord{-}\beta &\text{se }\alpha < 0_\Real\text{ e }\beta < 0_\Real\\
		\mathord{-}(\mathord{-}\alpha \times \beta) &\text{se }\alpha < 0_\Real \text{ e }\beta > 0_\Real\\
		\mathord{-}(\alpha \times \mathord{-}\beta) &\text{se }\alpha > 0_\Real \text{ e }\beta < 0_\Real
	\end{cases}
\end{align*}
Depois há que verificar que cada uma destas definições produz sempre um corte. E,
por fim, há que fazer uma demonstração fácil (mas longa) de que os cortes, assim
definidos, se comportam exatamente como devem. Mas relegamos tudo isto a
\olref[check]{sec}.
\end{document}
